% Part: first-order-logic
% Chapter: models-theories
% Section: expressing-props-of-structures

\documentclass[../../../include/open-logic-section]{subfiles}

\begin{document}

\olfileid{fol}{mat}{exs}

\olsection{Eigenschappen van \printtoken{P}{structure} uitdrukken}

\begin{explain}
Vaak willen we kunnen zeggen aan welke voorwaarden functies en
relaties moeten voldoen. Nog algemener willen we kunnen zeggen dat de
functies en relaties in een bepaalde structuur aan die voorwaarden
voldoen. Daar is een reden voor. Bij een taal hebben we meestal voor
ogen wat de !!{predicate}s moeten ``betekenen.'' We willen de
!!{structure}s waarin ze die bedoelde betekenis hebben kunnen
onderscheiden van de !!{structure}s waarin dat niet zo is.
Welke !!{structure}s we bedoelen, mogen we natuurlijk zelf kiezen. We
kunnen bijvoorbeeld zeggen dat de interpretatie van het
!!{predicate}~$\le$ een ordening moet zijn. Of we kunnen alleen kijken
naar interpretaties van~$\Lang L$ waarbij het domein uit verzamelingen
bestaat en $\Obj \in$ de relatie ``is !!a{element} van'' betekent. De
vraag is of we zo'n keuze ook met !!{sentence}s van de taal kunnen
vastleggen. Preciezer gezegd: welke voorwaarden aan
!!a{structure}~$\Struct M$ kunnen we uitdrukken met !!a{sentence} (of
misschien een verzameling !!{sentence}s) in de taal van~$\Struct M$?
Dat lukt niet voor elke voorwaarde. Er bestaat bijvoorbeeld geen zin
van~$\Lang L_A$ die in !!a{structure}~$\Struct M$ alleen waar is als
$\Domain M = \Nat$. We kunnen dus niet uitdrukken dat ``het domein
alleen natuurlijke getallen bevat.'' Sommige ``structurele
eigenschappen'' van !!{structure}s kunnen we misschien wel
uitdrukken. De vraag is welke eigenschappen van !!{structure}s we met
!!{sentence}s kunnen vastleggen. Anders gezegd: welke verzamelingen
!!{structure}s bestaan precies uit de structuren die !!a{sentence} (of
een verzameling !!{sentence}s) waar maken?
\end{explain}

\begin{defn}[Model van een verzameling]
Neem een verzameling !!{sentence}s~$\Gamma$ in een taal~$\Lang L$. We
noemen !!a{structure}~$\Struct M$ \emph{een model van}~$\Gamma$ als
$\Sat{M}{!A}$ voor alle $!A \in \Gamma$. Alle zinnen uit de verzameling
zijn dan dus waar in die structuur.
\end{defn}

\begin{ex}
De zin $\lforall[x][x \le x]$ is waar in~$\Struct M$ precies als
$\Assign{\le}{M}$ reflexief is: elk object staat dus in deze relatie
tot zichzelf. De zin
$\lforall[x][\lforall[y][((x \le y \land y \le x) \lif x = y)]]$ is
waar in~$\Struct M$ precies als $\Assign{\le}{M}$ antisymmetrisch is:
als twee objecten in beide richtingen in de relatie staan, zijn ze
hetzelfde. De zin
$\lforall[x][\lforall[y][\lforall[z][((x \le y \land y \le z)
      \lif x \le z)]]]$ is waar in~$\Struct M$ precies als
$\Assign{\le}{M}$ transitief is: als een eerste object in de relatie
staat tot een tweede, en dat tweede tot een derde, dan staat het eerste
ook in de relatie tot het derde. De modellen van
\begin{align*}
\{\quad &\lforall[x][x \le x], \\
   & \lforall[x][\lforall[y][((x \le y \land y \le
    x) \lif x = y)]], \\
   &\lforall[x][\lforall[y][\lforall[z][((x \le y
      \land y \le z) \lif x \le z)]]] \quad \}
\end{align*}
zijn dus precies de structuren waarin~$\Assign{\le}{M}$ reflexief,
antisymmetrisch en transitief is. Samen betekent dat dat de relatie een
partiële ordening is. We kunnen deze zinnen daarom als axioma's nemen
voor de \emph{eerste-ordetheorie van partiële ordeningen}.
\end{ex}

\end{document}
